% 张三慧《学习辅导与习题解答》扫描 S 卷1第8章“思考题选答”原编号。
% 仅辅导书来源已视觉核验，V1主教材原题逐图交叉核验仍待；全部exam_ready=false。
\section*{第8章\quad 狭义相对论基础：辅导书已列思考题独立解析}
\paragraph{思考题 8.3}
设列车参考系 \(S'\) 沿 \(+x\) 相对地面 \(S\) 以速度 \(v>0\) 运动，车头在 \(x'_{\rm front}>x'_{\rm rear}\)。由逆Lorentz变换 \(t=\gamma(t'+vx'/c^2)\)，若两闪电 \(\Delta t'=0\)，则 \(\boxed{\Delta t=\gamma v\,\Delta x'/c^2>0}\)，即按\emph{此坐标约定}地面系中车头（前端）闪电较晚、车尾较早。若列车运动方向反转则先后相反。需说明同时性的参照系，不能脱离列车正向和图像就报“车尾早”。

\paragraph{思考题 8.5}
本题应\emph{先指定S系的一组同时事件}，如 \(t=0\) 时所有S系钟 \(C_1,\ldots,C_5\) 同步，均指零。设S'沿S的正x方向运动，Lorentz变换 \(t'=\gamma(t-vx/c^2)\)。在S系 \(t=0\) 的同一时刻，位于 \(x'_j\) 处的运动钟读数为 \(\boxed{t'_j=-v x'_j/c^2}\)（因为 \(x=\gamma x'\) 在S系的同一时刻成立），故\emph{S'系靠左即 \(x'<0\) 的钟读数为正，靠右 \(x'>0\) 的钟读数为负}；居中 \(C_3'\) 读数为零。各钟在S'系自身当然仍然是同步的；不同系“同一时刻”的选取不同。若五钟的S'静止系间距为 \(d\)，则从左到右分别为 \(+2vd/c^2,+vd/c^2,0,-vd/c^2,-2vd/c^2\)。原图钟针仅作定性说明，未给出d、v的数值，不能编造具体秒数。

\paragraph{思考题 8.7}
在某惯性系里量杆长度，须记录\emph{同一时刻}两端坐标并取差 \(L=x_2(t)-x_1(t)\)。对于沿杆轴相对运动且静止固有长 \(L_0\) 的杆，\(\boxed{L=L_0/\gamma=L_0\sqrt{1-v^2/c^2}}\)。此是不同惯性系选取不同同时切片的时空测量结果，不需要弹性应力或材料压缩；钢杆与木杆有相同的理想运动学收缩因子。
